% 平面计算图采用浅色圆角模块、简短算子标签和连贯的圆角跨层路径。
\tikzset{
  nn operator/.style={draw=NNHidden,fill=NNHidden!18,rounded corners=1.5mm,
    line width=1.1pt,minimum height=10mm,inner sep=0pt,align=center},
  nn operator path/.style={nn flow,draw=NNHidden,rounded corners=2.5mm},
  nn pointwise/.style={circle,draw=NNHidden,fill=NNCanvas,line width=1.1pt,
    minimum size=6mm,inner sep=0pt}
}
% 名称、中心 x/y、核尺寸。
\newcommand{\nnConvOp}[4]{%
  \node[nn operator,fill=NNHidden!18,minimum width=10mm,minimum height=13mm]
    (#1) at (#2,#3) {Conv\\$#4$};%
}
\newcommand{\nnReluOp}[3]{%
  \node[nn operator,fill=NNHidden!18,minimum width=10mm] (#1) at (#2,#3) {ReLU};%
}
\newcommand{\nnSigmoidOp}[3]{%
  \node[nn operator,fill=NNHidden!18,minimum width=17mm] (#1) at (#2,#3) {Sigmoid};%
}
\newcommand{\nnBNOp}[3]{%
  \node[nn operator,fill=NNHidden!18,minimum width=8mm] (#1) at (#2,#3) {BN};%
}
% 可选圆直径、名称、中心 x/y；青色圆形 + 表示通道拼接，图注说明其语义。
\newcommand{\nnConcatOp}[4][6mm]{%
  \node[nn pointwise,draw=NNOutput,minimum size=#1]
    (#2) at (#3,#4) {$+$};%
}
\newcommand{\nnTensorState}[5]{%
  \node[draw=#4,fill=#4!8,line width=1.1pt,minimum width=8mm,
    minimum height=20mm,inner sep=0pt] (#1) at (#2,#3) {};%
  \foreach \nnOffset in {-6,0,6}{%
    \node[circle,draw=#4,fill=#4!20,line width=1.1pt,
      minimum size=3mm,inner sep=0pt] at (#2,{#3+\nnOffset}) {};%
  }%
  \node[nn heading,anchor=south] at (#2,{#3+12}) {#5};%
}
